\documentclass{article}
\usepackage{amsmath,amssymb,amsthm}
\newtheorem{theorem}{Theorem}
\title{Actual arena evaluation: local mapping theorem and missing contraction cells}
\date{16 November 2026}
\begin{document}\maketitle
\paragraph{Source and scope.}
This reads \texttt{Incremental\_Causal\_Arenas\_and\_Open\_Control.tex} at
SHA-256 \texttt{67ce4c52a14c740f0e9271d73980e0c494adae6dd93fdc02f27e9ee0d0f69c42}.
Its arenas have genuinely incremental output ports and open component
composition. Their native maps preserve enabled events and emitted ports
bijectively at every configuration. Write $\mathcal A(A,B)$ for these native
hom categories and $E:\mathcal A(A,B)\to H(A,B)$ for independent evaluation.
The results below use the claimed native evaluation functor, not a selected
shared-policy evaluation as a horizontal functor. Independent evaluation's
full proof is undergoing separate review. No effectful symmetric tensor is
assumed. This note separates a proved obstruction inside the supplied arena
category from a standard, precisely limited mapping construction.

\section{The actual native category has no contraction reflection}
\begin{theorem}
The atomic inclusion $I:H(A,B)\to\mathcal A(A,B)$ is not a right adjoint to
$E$ in general. On the supplied native arena-map class it does not even admit
general lax or oplax unit/composition comparison cells.
\end{theorem}
\begin{proof}
Use a singleton interface and the original identity process $1$, whose
atomic arena $I(1)$ has a singleton enabled event with singleton raw outcome,
followed by a terminal configuration with one emitted port. The actual arena
identity has no events. No native map in either direction can compare these
initial states, since their enabled sets have sizes zero and one.
Thus neither orientation of a unit comparison exists.

Let $R=I(1)\circ I(1)$. Its initial state has one enabled event. After that
event, the first component has emitted its port and the second component has
one enabled event. A native map $R\to I(1)$ would send the first singleton
event to the atomic event. Its transition equation would then compare the
intermediate state of $R$, which has one enabled event, to the terminal
atomic state, which has none. This contradicts the enabled-set bijection.
The same transition argument excludes a map $I(1)\to R$.
But $E(R)$ is canonically the original identity process, so
$H(E(R),1)$ contains its canonical native identity/unitor. Hence the required
adjunction bijection
$\mathcal A(R,I(1))\cong H(E(R),1)$ cannot exist. The missing comparison
cells also exclude general lax and oplax source compositors with the given
atomic assignment. This example uses only unary unit grades, deterministic
rows and actual transition squares; probability or infinite behavior is
irrelevant.
\end{proof}

The semantic retraction $EI\cong1$ still proves faithful source recovery;
it does not construct an operational contraction cell. Enlarging the map
class is a necessary new construction, not a coherence choice omitted from
the current definition. The stronger shared-scope obstruction in the
companion note shows why freely whiskering such added cells through all
controlled contexts is not semantically justified either.

\begin{theorem}[Unit-preserving context recovery is impossible in this map class]
Let $M(1,1)$ be the original normalized native context category, with its
canonical embedding into $H$ using unit content. There is no functor
$J:M(1,1)\to\mathcal A(1,1)$ simultaneously satisfying: $J$ sends the
context unit to an arena operationally isomorphic to the zero-event identity;
$J$ preserves all native context arrows; and $EJ$ is natively isomorphic to
the original context embedding. This holds independently of the atomic
choice of implementation.
\end{theorem}
\begin{proof}
Choose the context $P$ with two raw points of weights $1/2,1/2$, unit
coefficients and constant output. The actual native augmentation collapse
$P\to1$ exists. Its image is an arena map $J(P)\to J(1)$. At the designated
initial configuration it must induce a bijection of enabled sets. Since
$J(1)$ is operationally isomorphic to the zero-event identity, its initial
enabled set is empty; therefore so is the initial enabled set of $J(P)$.
That initial configuration is terminal. Independent evaluation consequently
has one raw witness, the empty execution trace, with weight one. Configurations
outside initial reach cannot alter this evaluated row. A native isomorphism
to the original two-point row $P$ would require a raw bijection, which is
impossible. The contradiction uses actual native recovery, not equality of
scalar output distributions.
\end{proof}
This pins the missing structure more strongly than the atomic counterexample:
retaining original native collapses, operational unit preservation and exact
native retraction cannot coexist in this enabled-bijection map doctrine.
Changing the implementation assigned to $P$ does not remove that obstruction.

\section{A genuine local universal property, at its stated dimension}
Fix endpoints $A,B$ and the actual functor $E$. Define the evaluation
collage $\mathcal C_E$ with two disjoint kinds of objects: operational arenas
$P\in\mathcal A(A,B)$ and opaque processes $[F]$ for $F\in H(A,B)$.
Its hom sets are
\[
\begin{array}{ll}
\mathcal C_E(P,Q)=\mathcal A(P,Q),&
\mathcal C_E([F],[G])=H(F,G),\\
\mathcal C_E(P,[F])=H(EP,F),&
\mathcal C_E([F],P)=\varnothing.
\end{array}
\]
Native composition is unchanged. For a mixed arrow $\alpha:EP\to F$,
precomposition by $u:Q\to P$ is $\alpha E(u)$, and postcomposition by
$\beta:F\to G$ is $\beta\alpha$. Functoriality of $E$ proves associativity.
Let $K:\mathcal A\to\mathcal C_E$ and $J:H\to\mathcal C_E$ be the two
inclusions. The canonical mixed arrow $\eta_P:P\to[EP]$ is $1_{EP}$.

\begin{theorem}[Standard evaluation-collage characterization]
Both $K$ and $J$ are fully faithful. Evaluation extends to
$\bar E:\mathcal C_E\to H$, taking $P$ to $EP$ and $[F]$ to $F$, and
$\bar E\dashv J$ with identity counit and unit $\eta$ on operational
objects. For every category $\mathcal D$, specifying a functor
$T:\mathcal C_E\to\mathcal D$ is equivalent to specifying functors
$F:\mathcal A\to\mathcal D$, $G:H\to\mathcal D$, and a natural
transformation $\theta:F\Rightarrow GE$.
This equivalence also identifies natural transformations: pairs of
transformations $(a:F\Rightarrow F',b:G\Rightarrow G')$ must satisfy
$(bE)\theta=\theta'a$.
\end{theorem}
\begin{proof}
Full faithfulness follows from the displayed hom sets. For every object $X$
of the collage, the defining equation
$\mathcal C_E(X,[F])=H(\bar EX,F)$ is natural in both variables and proves
the adjunction. Its unit and counit have the asserted descriptions.

Restriction gives $F=TK$, $G=TJ$, $\theta=T\eta$; naturality follows from
$\eta_Q u=(Eu)\eta_P$. Conversely send a mixed arrow $\alpha:EP\to F_0$
to $G(\alpha)\theta_P$. Naturality of $\theta$ makes this compatible with
precomposition, and functoriality of $G$ handles postcomposition. There are
no other mixed compositions. Every mixed arrow factors uniquely as
$J(\alpha)\eta_P$, so this extension is unique. The natural-transformation
statement follows by checking that same factorization.
\end{proof}

This is a concrete instance of the standard collage of a representable
profunctor, not a novel universal construction. It preserves all original
native maps and all operational native maps by full embeddings. Its opaque
$[F]$ is deliberately \emph{not} the existing operational atom $I(F)$:
there is a directed arrow $I(F)\to[F]$, but no reverse arrow. Thus the theorem
does not contradict the preceding obstruction or secretly identify those
interfaces. Full grades and their actual maps are retained by the supplied
$E$ and by the entire $H(EP,F)$, not reduced to terminal scalar laws.

\section{What this does not establish}
The collage theorem is local in $(A,B)$, at the category of native
comparisons. No component-substitution operation on these collages has been
constructed here. Defining every composition involving an opaque object to
be opaque would supply a different, absorbing evaluation convention; it would
not prove preservation of the original open intervention ports. Similarly,
a free bicategory with formal contraction cells and chosen equations needs
a separate semantic interpretation before any contextual compatibility is
claimed.

The intrinsic arena class may include persistent fork/join structures and
early-output behavior not proved generated by original binary $H$ composition.
Accordingly the collage takes the entire arena doctrine as an additional
input. Neither semantic retraction nor this mapping property is a free
completion theorem for $H$. Such a claim requires a separate generation-and-
relations theorem, or a declared independent event/port generator doctrine.

The actual arena bicategory already supplies independently characterized
open substitution and a strong independent evaluator, if its source proof
passes review. The local collage gives an exact relative mapping property
for adjoining whole-implementation semantic comparisons. Combining the two
into a contraction-compatible bicategory is an unresolved compatibility
obligation, not settled by their separate existence. The no-go examples rule
out unrestricted scope-reflecting comparison and unrestricted controlled
whiskering, while leaving sealed/retained-interface constructions open.
Selected policies also do not acquire horizontal composition merely because
their native-map compatibility relations compose. Intermediates in those
relations must share the same retained open scope; closing a component's
history can lose dependencies on another component's revealed outcomes.

\paragraph{Audit amendment for independent native transport.}
For a native arena map, the source proof should choose a deterministic target
scheduler and pull its selected subsets back through the enabled-event
bijections at each source history. Choosing a source scheduler first and
pushing it forward can conflict when two source histories merge. Pullback
is well-defined on every source history, including zero and unreachable
ones, because target tables are complete and event maps are bijections.
With independent laws, primitive exact pushforwards then prove the finite
terminal fibre equations. This is a local proof repair, not a restriction of
the admitted native maps.
\end{document}
